\documentclass[11pt]{article}
\usepackage[margin=1in]{geometry}
\usepackage{graphicx}
\usepackage{booktabs}
\usepackage{float}
\usepackage{hyperref}
\usepackage{amsmath}

\title{NHANES Adult Health Data: Race-Specific Diabetes Risk Profiles}
\author{Group Members: TODO}
\date{STAT 24620 / FINM 34700 / STAT 32950 -- Spring 2026}

\begin{document}
\maketitle

\section{Research Question}
TODO: State the final research question. Current tracked analysis asks whether cardiometabolic and socioeconomic diabetes risk profiles differ across self-identified White and Black individuals in the NHANES adult health data.

\section{Dataset and Preprocessing}
TODO: Describe \texttt{data/nhanes\_health.csv}, complete-case sample, duplicate removal from 2,412 raw rows to 1,508 unique profiles, and the no-survey-weights limitation.

\section{Descriptive Exploration}
TODO: Summarize sample composition, diabetes/high-BP rates, socioeconomic differences, and correlation structure. Select 1--2 figures/tables.

\section{Multivariate Structure}
TODO: Compare PCA, factor analysis, and k-modes clustering. Explain what each method answers and what it reveals.

\section{Supervised Comparison}
TODO: Summarize race-specific logistic regression for \texttt{Diabetes = Yes}, including AUC, sensitivity/specificity, and coefficient patterns.

\section{Method Comparison}
TODO: Explicitly compare PCA/FA/clustering/logistic regression: structure vs. latent factors vs. categorical profiles vs. prediction.

\section{Limitations}
TODO: Discuss cross-sectional data, no survey weights, sample-size imbalance, duplicate handling, and association-not-causation.

\end{document}
